% APA 7 (student) research paper
% Generated from Proposal.pdf scope on 2026-04-15

\documentclass[stu,man]{apa7}

\usepackage[utf8]{inputenc}
\usepackage{csquotes}
\usepackage{graphicx}
\usepackage{microtype}
\usepackage{booktabs}
\usepackage{hyperref}

% Use BibLaTeX for APA references
\usepackage[style=apa,backend=biber,sortcites=true,sorting=nyt]{biblatex}
\DeclareLanguageMapping{english}{english-apa}
\addbibresource{references.bib}

\title{Outfit Recommendation System Based on Weather: \texorpdfstring{\\}{ }An AI-Assisted Virtual Wardrobe with Rule-Based Comfort Scoring}

\authorsnames{Talha Mushtaq \and Fizza Ahmad \and Muhammad Shabih Haider \and Shahzaib Sajid}
\authorsaffiliations{Department of Computer Science, University of Engineering and Technology Lahore}

\authornote{Supervisor: Ms. Darakhshan Bokhat\\
Session: 2022--2026}

\abstract{Choosing what to wear is a daily decision that becomes harder when weather conditions vary and when users lack explicit styling knowledge. This paper presents an AI-assisted web application that recommends outfits using the user\textquotesingle s own wardrobe, real-time weather data, and foundational fashion principles.
The system combines (a) wardrobe management for user-uploaded clothing images, (b) visual attribute extraction and color analysis, and (c) an explainable rule-based scoring layer that operationalizes comfort and seasonal guidelines (e.g., thermal and water-vapour resistance concepts from ISO testing practice) to produce climate-appropriate, aesthetically coherent combinations.
Unlike generic fashion recommenders that ignore ownership constraints, the proposed approach generates recommendations grounded in the clothing items the user actually owns, improving practicality and user trust.
We describe the end-to-end methodology, implementation choices, and validation evidence from automated tests of the rule-based engine and weather scoring components.}

\keywords{fashion recommendation, weather-aware recommendation, virtual wardrobe, rule-based scoring, computer vision}

\begin{document}
\maketitle

\section{Introduction}
Daily outfit selection is often treated as a simple routine, yet it becomes a recurring decision problem under shifting weather conditions and diverse social contexts.
Many existing applications provide generic outfit templates or rely on catalog inventories rather than personal wardrobes, which can make recommendations impractical.
A more useful recommender must (1) be constrained by what the user actually owns, (2) adapt to real-time weather inputs, and (3) provide intelligible reasoning so users can trust and refine suggestions.

This research proposes an AI-assisted outfit recommendation system that integrates user wardrobe images with live weather data and a transparent scoring layer.
The proposal is aligned with modern computer vision foundations for attribute learning \parencite{NIPS2012_c399862d} and leverages established fashion datasets such as DeepFashion \parencite{liuLQWTcvpr16DeepFashion} and DeepFashion-MultiModal \parencite{deepfashionmultimodal}.
The resulting system targets daily use: recommending outfits optimized for comfort, practicality, and visual coherence.

\section{Related Work}
Modern fashion recommendation systems employ a variety of AI techniques to improve personalization and context-awareness \parencite{fashion2023review}.
A key challenge in outfit recommendation is the integration of visual perception with contextual reasoning (weather, occasion, personal preference).

Large-scale labeled datasets enable learning clothing categories and attributes from images.
DeepFashion provides rich annotations (categories, attributes, bounding boxes, and landmarks) supporting recognition and retrieval tasks \parencite{liuLQWTcvpr16DeepFashion}.
DeepFashion-MultiModal extends fashion understanding with multi-modal labels such as parsing masks, keypoints, and clothing texture attributes \parencite{deepfashionmultimodal,jiang2022text2human}.

Weather-aware recommendations depend on reliable meteorological data sources.
In practice, production systems often integrate commercial weather APIs that provide current conditions and forecasts via standardized JSON responses \parencite{openweatherapi}.
Existing commercial applications (e.g., outfit planning apps and wardrobe management systems) offer varying levels of weather awareness, but many lack transparent reasoning or scientific grounding in comfort standards.

From a systems perspective, training and deploying deep models frequently relies on frameworks such as TensorFlow \parencite{abadi2016tensorflow}.
For classical machine learning utilities (e.g., clustering and baseline models), scikit-learn provides a widely used, well-documented implementation ecosystem \parencite{pedregosa2011scikit}.
Finally, clustering of mixed numeric and categorical features (e.g., weather variables and occasion labels) can be addressed by extensions to $k$-means that explicitly handle categorical values \parencite{huang1998extensions}.

\section{Method}
\subsection{Participants}
This work is a system design and engineering study.
No human-subject experiment was conducted for the current submission.
The system is intended for end users managing personal wardrobes; formal user studies are planned as future work.

\subsection{Materials}
\textbf{Wardrobe data.} Users provide images of clothing items (e.g., tops, bottoms, outerwear). Images are stored as part of a virtual wardrobe and associated with inferred attributes (color, type, and comfort-related properties).

\textbf{Training and reference data.} For model development and benchmarking, we reference publicly available fashion datasets such as DeepFashion \parencite{liuLQWTcvpr16DeepFashion} and DeepFashion-MultiModal \parencite{deepfashionmultimodal}.

\textbf{Weather data.} Real-time meteorological variables (temperature, humidity, and conditions such as rain) are fetched from OpenWeather APIs \parencite{openweatherapi}.

\textbf{Software stack.} The system is implemented in Python with a web backend and an ML pipeline.
TensorFlow supports deep learning workflows \parencite{abadi2016tensorflow}; OpenCV is used for image preprocessing utilities \parencite{opencvdocs}; and scikit-learn supports classical ML components \parencite{pedregosa2011scikit}.

\subsection{Design}
The proposed architecture is a hybrid of learned perception and explainable decision rules:
\begin{enumerate}
  \item \textbf{Perception layer}: attribute extraction from images (e.g., garment category, sleeve length, and visual features as proxies for fabric type inference).
  \item \textbf{Context layer}: real-time weather features (temperature, humidity, precipitation, wind).
  \item \textbf{Rule-based scoring layer}: interpretable rules combining color harmony constraints and comfort heuristics.
  \item \textbf{Outfit generation layer}: constrained combinatorial search over the user\textquotesingle s wardrobe to produce outfits (top + bottom + optional outerwear) ranked by total score.
\end{enumerate}

The scoring layer is designed to be explainable: each recommendation can be decomposed into contextual sub-scores (e.g., warmth suitability, breathability suitability, rain appropriateness, color compatibility, and occasion appropriateness such as casual versus formal settings).
Comfort evaluation is inspired by textile testing concepts for thermal and evaporative resistance \parencite{iso11092}, but implemented as a practical rule engine rather than a laboratory measurement instrument.

\section{Procedure}
\subsection{Requirement Gathering}
We derive core requirements from the project proposal: (1) personalized recommendations based on owned wardrobe items, (2) weather-aware adaptation, (3) color-theory matching, and (4) scientifically informed comfort rules.

\subsection{Data Collection and Preparation}
Public datasets (e.g., DeepFashion) provide labeled images for attribute learning and validation \parencite{liuLQWTcvpr16DeepFashion}.
For user wardrobes, images are normalized through a preprocessing pipeline (resizing, background normalization when available, and color extraction).

\subsection{Attribute Extraction}
A convolutional neural network (CNN) is used to infer visual attributes.
CNN-based representation learning is motivated by established results in large-scale image classification \parencite{NIPS2012_c399862d}.

\subsection{Weather Integration}
Weather data are retrieved on demand and transformed into a feature vector used by the scoring layer (e.g., mapping precipitation to rain-protection requirements).

\subsection{Outfit Scoring and Ranking}
Candidate outfits are enumerated from wardrobe categories.
Each candidate receives a total score:
\[
S = w_c S_{color} + w_w S_{weather} + w_k S_{comfort} + w_o S_{occasion},
\]
where $w_c, w_w, w_k, w_o$ denote configurable weights for color harmony, weather suitability, comfort, and occasion appropriateness respectively, and the corresponding $S$ components are normalized to $[0, 1]$.
The system returns the top-ranked outfits and the explanation trace of contributing factors.

\section{Results}
Because the current submission focuses on an engineering prototype, we report verifiable software validation results rather than user-study outcomes.

\subsection{Automated Tests}
A focused subset of unit tests covering weather scoring, color scoring, and rule-based constraints was executed under Python 3.11.
All tests in the selected suite passed (43/43) in a single run (0.78 seconds).
This provides evidence that core logic remains consistent across representative inputs and edge cases.

\subsection{Functional Output}
Given a populated wardrobe and a weather snapshot, the system produces a ranked list of outfits.
Each recommendation includes: (1) the selected garment items, (2) total score, and (3) a breakdown of weather suitability and comfort logic.

\section{Discussion}
The proposed system directly addresses common limitations of generic outfit recommenders by constraining suggestions to the user\textquotesingle s owned wardrobe and by incorporating live weather context \parencite{openweatherapi}.
The explainable scoring layer improves transparency by exposing reasons for each recommendation (e.g., “layering suggested due to low temperature” or “avoid absorbent fabrics during rain”).

\subsection{Limitations}
This work has several practical constraints:
\begin{itemize}
  \item \textbf{Fabric inference}: fabric type is difficult to infer reliably from images alone; learned predictions and rule assumptions can be uncertain.
  \item \textbf{Dataset bias}: public fashion datasets may not represent local clothing styles or cultural norms; this can affect generalization.
  \item \textbf{Weather granularity}: microclimates and indoor/outdoor transitions are not fully captured by a single API call.
\end{itemize}

\subsection{Future Work}
Future improvements include: (1) user-feedback loops for personalization, (2) occasion-aware context with calendar integration, (3) virtual try-on integration for confidence before wearing, and (4) rigorous user studies measuring decision time reduction and perceived recommendation quality.

\section{Conclusion}
This paper presents a weather-aware outfit recommendation system that combines wardrobe-constrained outfit generation with an explainable, rule-based comfort and style scoring layer.
Grounded in established datasets and modern ML tooling \parencite{liuLQWTcvpr16DeepFashion,abadi2016tensorflow,pedregosa2011scikit}, the proposed approach prioritizes practicality, transparency, and context awareness.
The system\textquotesingle s core scoring logic is validated through automated tests of the rule and weather components.

\printbibliography

\end{document}
